\documentclass[a4paper,10pt,titlepage,oneside]{article}

% ------------------------------------------------------------------
% METADATA
% ------------------------------------------------------------------
\title{}
\author{}
\date{\today}

% ------------------------------------------------------------------
% BASIC PACKAGES AND STYLE
% ------------------------------------------------------------------
\usepackage[utf8]{inputenc}
\usepackage[T1]{fontenc}
\usepackage[english]{babel}
\usepackage{tgpagella}
\usepackage{float}
\usepackage[a4paper,top=3cm,bottom=3.2cm,inner=3cm,outer=2cm,heightrounded]{geometry}
\usepackage{amsmath,amssymb,mathtools,graphicx,booktabs,fancyhdr,enumitem}
\usepackage[hidelinks]{hyperref}
\usepackage{xcolor}
\usepackage{longtable}
\usepackage{array}
\usepackage{listings}
\usepackage{titlesec}

\setlength{\parindent}{15pt}
\setlength{\parskip}{0pt}
\setlength{\footskip}{60pt}
\setlength{\headheight}{15pt}

\hypersetup{
    pdftitle={SOVD Raspberry Demo Technical Documentation},
    pdfauthor={Ana Lozano Castillo}
}

\definecolor{codebg}{RGB}{246,248,250}
\definecolor{codeframe}{RGB}{208,215,222}
\definecolor{headingblue}{RGB}{31,78,121}

\lstdefinestyle{projectcode}{
    backgroundcolor=\color{codebg},
    frame=single,
    rulecolor=\color{codeframe},
    basicstyle=\ttfamily\small,
    breaklines=true,
    columns=fullflexible,
    keepspaces=true,
    showstringspaces=false
}

\lstset{style=projectcode}

\titleformat{\section}
    {\Large\bfseries\color{headingblue}}
    {\thesection}{0.8em}{}

\titleformat{\subsection}
    {\large\bfseries}
    {\thesubsection}{0.8em}{}

% ------------------------------------------------------------------
% HEADERS AND FOOTERS
% ------------------------------------------------------------------
\fancyhf{}
\fancyhead[OL]{SOVD Raspberry Demo}
\fancyhead[OR]{Technical Documentation}
\fancyfoot[OR]{p. \thepage}
\pagestyle{fancy}
\raggedbottom

\renewcommand{\thesection}{\arabic{section}}
\renewcommand{\thesubsection}{\thesection.\arabic{subsection}}

% ------------------------------------------------------------------
% DOCUMENT
% ------------------------------------------------------------------
\begin{document}

% ------------------------------------------------------------------
% TITLE PAGE
% ------------------------------------------------------------------
\begin{titlepage}
\vspace*{1.5cm}
{\centering
{\Huge \textbf{SOVD Raspberry Demo}}\\[8mm]
{\Large \textbf{Technical Project Documentation}}\\[1.5cm]
{\large Service-Oriented Vehicle Diagnostics prototype for Raspberry Pi}\\[2cm]
\begin{tabular}{l@{\hspace{15mm}}l}
\Large \textbf{Author:} & \Large Ana Lozano Castillo\\[5mm]
\Large \textbf{Date:} & \Large \today\\
\end{tabular}\\[3cm]
\vfill
{\large Flask REST API \quad DoIP ECU Simulator \quad UDS Diagnostics \quad Arduino Bench}\\[8mm]
}
\end{titlepage}

% ------------------------------------------------------------------
% TABLE OF CONTENTS
% ------------------------------------------------------------------
\pagenumbering{roman}
\renewcommand{\contentsname}{Table of Contents}
\tableofcontents
\newpage

\pagenumbering{arabic}
\setcounter{page}{1}

\section{Executive Summary}

The SOVD Raspberry Demo is an educational vehicle diagnostics project designed
for Raspberry Pi based experimentation. It combines a Flask web dashboard, a
REST API, a simulated DoIP ECU, UDS diagnostic services, real-time browser
events, and an Arduino-based hardware bench.

The project demonstrates how service-oriented diagnostics can be represented in
a compact prototype. A user can open the dashboard, inspect simulated vehicle
data, execute UDS \texttt{ReadDataByIdentifier} requests through DoIP, and
receive live notifications when physical bench events occur, such as LED
disconnects, crash/touch input, or battery voltage faults.

The system is intentionally lightweight: state is held in memory, the ECU is
simulated in Python, and the dashboard is served by a Flask application. This
makes the project suitable for demonstrations, coursework, prototyping, and
hands-on diagnostics training.

\section{Project Objectives}

The main objectives of the project are:

\begin{itemize}[leftmargin=*]
    \item Provide a browser-based diagnostics dashboard on port \texttt{5000}.
    \item Expose REST endpoints for vehicle status, power data, components, and simulation controls.
    \item Simulate an ECU reachable through DoIP on TCP port \texttt{13400}.
    \item Support selected UDS \texttt{ReadDataByIdentifier} and \texttt{WriteDataByIdentifier} flows.
    \item Integrate an Arduino bench for physical LED, crash, and battery voltage events.
    \item Push real-time notifications to the browser through Server-Sent Events.
    \item Keep the dashboard and DoIP ECU running automatically through \texttt{systemd} services.
\end{itemize}

\section{System Architecture}

The project is organized around four main runtime components:

\begin{enumerate}[leftmargin=*]
    \item \textbf{Flask Dashboard Server}: implemented in \texttt{server.py}. It serves the HTML dashboard, REST API, UDS endpoint, Arduino serial listener, and SSE stream.
    \item \textbf{Simulated DoIP ECU}: implemented in \texttt{doip\_ecu.py}. It listens on TCP port \texttt{13400} and responds to UDS diagnostic messages.
    \item \textbf{DoIP Client Helper}: implemented in \texttt{doip\_client.py}. It is used by the REST layer to send UDS payloads to the ECU simulator.
    \item \textbf{Arduino Bench Monitor}: implemented in \texttt{arduino\_code.cpp}. It monitors physical inputs and sends serial messages to the Flask server.
\end{enumerate}

\subsection{High-Level Data Flow}

A typical diagnostic flow follows these steps:

\begin{enumerate}[leftmargin=*]
    \item The user selects a DID in the dashboard and clicks \texttt{Execute ReadDID}.
    \item The browser sends a \texttt{POST /uds/readDataByIdentifier} request to the Flask server.
    \item The Flask route delegates the request to the API routing layer.
    \item \texttt{doip\_client.py} builds a DoIP diagnostic message containing the UDS payload.
    \item The payload is sent to the simulated ECU at \texttt{127.0.0.1:13400}.
    \item \texttt{doip\_ecu.py} processes the UDS request and returns a positive or negative response.
    \item The dashboard updates its diagnostic console with HTTP, UDS, and DoIP trace entries.
\end{enumerate}

\section{Repository Structure}

\begin{longtable}{>{\ttfamily}p{0.34\linewidth} p{0.58\linewidth}}
\textbf{Path} & \textbf{Purpose} \\
\hline
\endfirsthead
\textbf{Path} & \textbf{Purpose} \\
\hline
\endhead
server.py & Main Flask application. Serves the dashboard, API, UDS endpoint, serial monitor, test endpoints, and SSE stream. \\
index.html & Browser dashboard for vehicle identity, telemetry, diagnostic traces, UDS reads, simulation controls, and fault popups. \\
doip\_ecu.py & Simulated DoIP ECU listening on TCP port \texttt{13400}. Handles UDS read/write services. \\
doip\_client.py & TCP DoIP helper used by the Flask API to send UDS requests to \texttt{127.0.0.1:13400}. \\
arduino\_code.cpp & Arduino sketch for monitoring rear LED, front LED, crash/touch input, and battery voltage. \\
routes/api\_routes.py & API routing logic for vehicle data, power data, components, simulation controls, and UDS reads. \\
routes/ui\_routes.py & UI route mapping for \texttt{/} and \texttt{/ui}. \\
data/simulated\_data.py & In-memory vehicle, powertrain, component, sensor, and fault data. \\
data/diagnostic\_trace.py & Bounded in-memory diagnostic trace buffer using \texttt{deque(maxlen=100)}. \\
ops/run\_forever.sh & Local watchdog script that restarts \texttt{server.py} if it exits. \\
ops/sovd-dashboard.service & \texttt{systemd} unit for running the dashboard service at boot. \\
ops/sovd-doip-ecu.service & \texttt{systemd} unit for running the simulated DoIP ECU at boot. \\
README.md & General project guide, execution instructions, endpoints, and verification commands. \\
Dockerfile & Container definition included with the repository. \\
project\_documentation.tex & This technical documentation source file. \\
\end{longtable}

\section{Main Software Components}

\subsection{\texttt{server.py}: Flask Dashboard and Gateway}

\texttt{server.py} is the main application process. It acts as both the user
interface backend and the diagnostics gateway. Its responsibilities include:

\begin{itemize}[leftmargin=*]
    \item Binding the Flask application to \texttt{0.0.0.0:5000}.
    \item Serving the browser dashboard.
    \item Providing REST endpoints for simulated vehicle data.
    \item Receiving UDS requests from the dashboard.
    \item Calling the DoIP client helper to reach the simulated ECU.
    \item Listening for Arduino serial messages.
    \item Publishing real-time fault events through \texttt{/events}.
    \item Maintaining consistent fault state for LEDs, crash input, and battery voltage.
\end{itemize}

The application stores runtime state in Python module-level dictionaries and
queues. There is no database, which keeps the system simple and appropriate for
demonstration use.

\subsection{\texttt{doip\_ecu.py}: Simulated ECU}

\texttt{doip\_ecu.py} simulates a DoIP-capable ECU. It is intentionally kept
separate from the Flask server so it does not share responsibilities with the
web layer or compete for the Arduino serial port.

The ECU supports the following UDS services:

\begin{itemize}[leftmargin=*]
    \item \texttt{0x22}: \texttt{ReadDataByIdentifier}.
    \item \texttt{0x2E}: \texttt{WriteDataByIdentifier}.
    \item \texttt{0x3E}: \texttt{TesterPresent}.
\end{itemize}

\subsection{\texttt{doip\_client.py}: Diagnostic Transport Helper}

\texttt{doip\_client.py} builds DoIP diagnostic messages and sends them over
TCP to the ECU simulator. The configured destination is:

\begin{lstlisting}
127.0.0.1:13400
\end{lstlisting}

For example, reading DID \texttt{F190} sends the UDS payload:

\begin{lstlisting}
22F190
\end{lstlisting}

\subsection{\texttt{index.html}: Browser Dashboard}

The dashboard provides the user-facing diagnostic experience. It includes:

\begin{itemize}[leftmargin=*]
    \item Vehicle identity and status information.
    \item Power and bus telemetry.
    \item A DID selector and \texttt{Execute ReadDID} action.
    \item A terminal-style diagnostic trace view.
    \item Simulation controls for ignition and diagnostic state.
    \item Real-time fault popups.
    \item An SSE connection to \texttt{/events}.
\end{itemize}

\section{Arduino Bench Integration}

The Arduino sketch represents the physical bench. It monitors hardware inputs
and sends text messages over serial at \texttt{9600} baud.

\begin{longtable}{>{\ttfamily}p{0.18\linewidth} p{0.72\linewidth}}
\textbf{Pin} & \textbf{Function} \\
\hline
A0 & Rear LED connection state. \\
A2 & Front LED connection state. \\
D2 & Crash/touch input. \\
A1 & Battery voltage potentiometer. \\
\end{longtable}

\texttt{server.py} scans the following serial devices:

\begin{lstlisting}
/dev/ttyACM0
/dev/ttyACM1
/dev/ttyUSB0
/dev/ttyUSB1
\end{lstlisting}

Supported serial messages include:

\begin{lstlisting}
LED_REAR:FAULT
LED_REAR:OK
LED_FRONT:FAULT
LED_FRONT:OK
CRASH:FAULT:<duration_ms>
BATTERY:VOLTAGE:<voltage>
BATTERY:UNDERVOLTAGE:<voltage>
BATTERY:OVERVOLTAGE:<voltage>
BATTERY:OK:<voltage>
\end{lstlisting}

\section{REST API}

\subsection{Vehicle and Power Data}

\begin{lstlisting}
GET /vehicle
GET /vehicle/power
\end{lstlisting}

These endpoints return simulated vehicle identity, ignition state, uptime,
ambient temperature, battery voltage, current, and power mode.

\subsection{Component Data}

\begin{lstlisting}
GET /components
GET /components/engine/ident
GET /components/engine/runtime
GET /components/engine/sensors
GET /components/engine/faults
\end{lstlisting}

These endpoints expose simulated component, engine, runtime, sensor, and fault
information.

\subsection{Simulation Controls}

\begin{lstlisting}
GET /ignition/on
GET /ignition/off
GET /faults/on
GET /faults/off
GET /vehicle/speed
GET /vehicle/speed/<value>
\end{lstlisting}

Browser-driven simulation state is scoped per browser client. The dashboard
stores a client identifier in \texttt{localStorage} and sends it as
\texttt{X-SOVD-Client}. This prevents actions from one browser from unexpectedly
changing another browser's simulated state.

\subsection{UDS ReadDataByIdentifier}

\begin{lstlisting}
POST /uds/readDataByIdentifier
Content-Type: application/json

{"did": "F190"}
\end{lstlisting}

This endpoint converts a REST request into a UDS request transported over DoIP.

\section{Supported Diagnostic Identifiers}

\begin{longtable}{>{\ttfamily}p{0.18\linewidth} p{0.72\linewidth}}
\textbf{DID} & \textbf{Description} \\
\hline
F190 & Vehicle Identification Number (VIN). \\
F187 & ECU software version. \\
F18C & ECU hardware version. \\
F40C & Current engine RPM. \\
F40D & Engine load value. \\
F40E & Engine coolant temperature. \\
F1A1 & Latest rear LED state. \\
F1A2 & Latest front LED state. \\
F1A3 & Latest crash state. \\
F1A4 & Latest battery state. \\
\end{longtable}

Bench state changes are also written back to the simulated ECU with UDS
\texttt{WriteDataByIdentifier}. For example, a rear LED state update uses DID
\texttt{F1A1} and a UDS request beginning with:

\begin{lstlisting}
2E F1 A1 ...
\end{lstlisting}

A positive response from the ECU begins with:

\begin{lstlisting}
6E F1 A1
\end{lstlisting}

\section{Real-Time Event Handling}

The dashboard subscribes to:

\begin{lstlisting}
GET /events
\end{lstlisting}

This endpoint uses Server-Sent Events and emits messages such as:

\begin{lstlisting}
data: LED_REAR:FAULT
data: LED_FRONT:FAULT
data: BATTERY:UNDERVOLTAGE:10.8
data: BATTERY:OVERVOLTAGE:14.8
data: CRASH:FAULT:66
data: SYSTEM:OK
\end{lstlisting}

When a browser reconnects or refreshes the page, the server immediately replays
active physical bench faults. This prevents the UI from showing a restored state
while a physical fault is still present.

\section{Fault Model}

\subsection{LED Faults}

The Arduino can report rear and front LED transitions:

\begin{lstlisting}
LED_REAR:FAULT
LED_REAR:OK
LED_FRONT:FAULT
LED_FRONT:OK
\end{lstlisting}

Fault messages trigger a dashboard popup and activate the fault background.
Restoration messages only produce a global restored state when all active
physical bench faults have cleared.

\subsection{Battery Voltage}

Battery status is voltage-driven:

\begin{lstlisting}
11.0 V <= voltage <= 14.5 V
\end{lstlisting}

\begin{itemize}[leftmargin=*]
    \item Below \texttt{11.0 V}: \texttt{UNDERVOLTAGE}.
    \item Above \texttt{14.5 V}: \texttt{OVERVOLTAGE}.
    \item Within range: \texttt{NORMAL}.
\end{itemize}

The same thresholds must remain aligned in \texttt{arduino\_code.cpp},
\texttt{server.py}, and \texttt{index.html}.

\subsection{Crash Event}

Crash or touch input events use the following message format:

\begin{lstlisting}
CRASH:FAULT:<duration_ms>
\end{lstlisting}

The dashboard shows a crash notification and marks the system as faulted until
the state is restored.

\section{Diagnostic Trace Model}

Diagnostic trace entries are stored in \texttt{data/diagnostic\_trace.py} using
a bounded \texttt{deque(maxlen=100)}. The dashboard renders these entries in a
terminal-style console.

Browser-initiated diagnostic activity is scoped by client ID, while physical
bench events are global. This means one user's UDS read does not pollute another
user's diagnostic console, but real bench events are visible to every connected
browser.

The \texttt{Clean Console} button is client-local. It stores a trace cursor in
\texttt{localStorage} and hides older entries for the current browser without
deleting shared trace data.

\section{Operation}

\subsection{Manual Startup}

Start the simulated DoIP ECU:

\begin{lstlisting}
python3 doip_ecu.py
\end{lstlisting}

Start the dashboard server:

\begin{lstlisting}
python3 server.py
\end{lstlisting}

Open the dashboard through the public ngrok URL:

\begin{lstlisting}
https://684ad3b98246.ngrok.app/
\end{lstlisting}

ngrok is the documented browser access method for the demo. This avoids
depending on local network routing, changing Raspberry Pi IP addresses, or guest
Wi-Fi restrictions.

\subsection{Persistent Services}

The repository includes \texttt{systemd} unit files so the dashboard and ECU can
start automatically after boot.

Install and enable the dashboard service:

\begin{lstlisting}
sudo cp ops/sovd-dashboard.service /etc/systemd/system/sovd-dashboard.service
sudo systemctl daemon-reload
sudo systemctl enable --now sovd-dashboard.service
sudo systemctl status sovd-dashboard.service
\end{lstlisting}

Install and enable the DoIP ECU service:

\begin{lstlisting}
sudo cp ops/sovd-doip-ecu.service /etc/systemd/system/sovd-doip-ecu.service
sudo systemctl daemon-reload
sudo systemctl enable --now sovd-doip-ecu.service
sudo systemctl status sovd-doip-ecu.service
\end{lstlisting}

Service logs are available through the system journal:

\begin{lstlisting}
journalctl -u sovd-dashboard.service
journalctl -u sovd-doip-ecu.service
\end{lstlisting}

\section{ngrok Access}

For presentation scenarios, the project uses ngrok to expose the dashboard
without relying on local network routing. The public dashboard URL is:

\begin{lstlisting}
https://684ad3b98246.ngrok.app/
\end{lstlisting}

This approach is useful when the Raspberry Pi is connected to a guest Wi-Fi
network or any environment that blocks direct communication between devices.

\section{Testing Without Arduino Hardware}

LED fault notifications can be tested without the physical Arduino bench:

\begin{lstlisting}
https://684ad3b98246.ngrok.app/test/fault
https://684ad3b98246.ngrok.app/test/ok
\end{lstlisting}

\section{Verification Procedures}

\subsection{Python Syntax Check}

\begin{lstlisting}
find . -maxdepth 3 -type f -name '*.py' -print0 | xargs -0 -n1 python3 -m py_compile
\end{lstlisting}

\subsection{Basic Endpoint Checks}

\begin{lstlisting}
curl https://684ad3b98246.ngrok.app/vehicle
curl https://684ad3b98246.ngrok.app/test/fault
curl https://684ad3b98246.ngrok.app/test/ok
\end{lstlisting}

\subsection{SSE Stream Check}

\begin{lstlisting}
curl --max-time 2 https://684ad3b98246.ngrok.app/events
\end{lstlisting}

\subsection{ReadDID Check}

\begin{lstlisting}
curl -X POST https://684ad3b98246.ngrok.app/uds/readDataByIdentifier \
  -H 'Content-Type: application/json' \
  -d '{"did":"F190"}'
\end{lstlisting}

A successful response for \texttt{F190} returns the simulated VIN:

\begin{lstlisting}
WVWZZZ12345678901
\end{lstlisting}

\section{Known Limitations}

\begin{itemize}[leftmargin=*]
    \item The Flask server and DoIP ECU are demonstration services, not production-grade deployments.
    \item Runtime state is held in memory and resets when the processes restart.
    \item \texttt{server.py} should be the only process that opens the Arduino serial port.
    \item \texttt{doip\_ecu.py} should remain focused on DoIP/UDS behavior and should not take over Flask or serial responsibilities.
    \item If \texttt{Execute ReadDID} reports \texttt{Connection refused}, the DoIP ECU service is probably not running.
    \item Browser access is documented through the public ngrok URL to avoid network isolation issues.
    \item The UI should not rely on a single \texttt{SYSTEM:OK} event to clear fault state; current telemetry must also be considered.
\end{itemize}

\section{Conclusion}

The SOVD Raspberry Demo provides a compact, end-to-end representation of a
vehicle diagnostics workflow. It links a web dashboard, REST API, DoIP transport,
UDS diagnostic services, simulated ECU behavior, Arduino hardware events, and
real-time browser notifications.

The project is intentionally small, but it captures the core relationships among
SOVD-style service access, DoIP transport, UDS payloads, and physical vehicle
signals. This makes it a practical platform for learning, demonstrations, and
future extensions.

\end{document}
